% ============================================================================
% Chapter 6: Introduction to Operating Systems
% ============================================================================

\chapter{Introduction to Operating Systems}
\label{ch:os-intro}

\begin{objectives}
  \item Define what an operating system is and why it is needed
  \item Explain the key functions of an operating system
  \item Name popular operating systems used today
  \item Understand the relationship between hardware, OS, and applications
\end{objectives}

\bytetip[tip]{An operating system is like the \textbf{manager} of a restaurant. You (the customer) don't talk to the chef (hardware) directly --- the manager (OS) takes your order and makes sure everything runs smoothly!}

% ---------------------------------------------------------------
\section{What Is an Operating System?}
\label{sec:what-is-os}
\index{operating system!definition}

An \hl{operating system} (OS) is the most important software on a computer. It manages all the hardware and software, and provides the interface you see and interact with (like the desktop, icons, and windows).

\begin{theorybox}[Definition: Operating System]
An \textbf{Operating System (OS)}\index{operating system} is system software that:
\begin{itemize}[leftmargin=*, label={\textcolor{ModuleBlue}{\faAngleRight}}, itemsep=2pt]
  \item Manages the computer's \textbf{hardware} (CPU, memory, storage, peripherals)
  \item Provides a \textbf{user interface} (desktop, menus, windows)
  \item Runs and manages \textbf{application software}
  \item Handles \textbf{files and folders} on your storage drives
  \item Provides \textbf{security} (user accounts, passwords, firewall)
\end{itemize}
\end{theorybox}

% --- Figure 6.1: OS as a Bridge ---
\begin{figure}[H]
\centering
\begin{tikzpicture}
  % Bottom layer: Hardware
  \node[draw=NavyBlue, fill=NavyBlue!15, rounded corners=4pt,
        minimum width=12cm, minimum height=1.5cm, line width=1.5pt,
        font=\small\bfseries\sffamily, text=NavyBlue] (hw) 
    {\faMicrochip\ Hardware --- CPU, RAM, Storage, Peripherals};
  
  % Middle layer: Operating System
  \node[draw=ModuleBlue, fill=ModuleBlue!15, rounded corners=4pt,
        minimum width=13cm, minimum height=1.8cm, line width=2pt,
        font=\bfseries\sffamily, text=ModuleBlue, above=0.8cm of hw] (os) 
    {\faWindows\ Operating System --- The Bridge};
  
  % Bridge arch decoration
  \foreach \x in {-4,-2,0,2,4}{
    \draw[ModuleBlue!30, line width=1pt] 
      ([xshift=\x cm, yshift=-0.1cm]os.south) arc (180:0:0.5cm and 0.3cm);
  }
  
  % Top layer: Application Software
  \node[draw=AccentTeal, fill=AccentTeal!12, rounded corners=4pt,
        minimum width=11cm, minimum height=1.5cm, line width=1.5pt,
        font=\small\bfseries\sffamily, text=AccentTeal, above=0.8cm of os] (apps) 
    {\faCode\ Application Software --- Word, Chrome, Games};
  
  % Arrows
  \draw[ModuleBlue, line width=2pt, <->] (hw.north) -- (os.south)
    node[midway, right, font=\scriptsize\sffamily, text=NavyBlue] {manages};
  \draw[ModuleBlue, line width=2pt, <->] (os.north) -- (apps.south)
    node[midway, right, font=\scriptsize\sffamily, text=NavyBlue] {supports};
  
  % Label
  \node[below=0.3cm of hw, font=\scriptsize\sffamily\itshape, text=NavyBlue] 
    {The OS sits between hardware and apps, managing communication between them.};
\end{tikzpicture}
\caption{The operating system acts as a bridge between hardware and applications}
\label{fig:os-bridge}
\end{figure}

% ---------------------------------------------------------------
\section{A Brief History of Windows}
\label{sec:windows-history}
\index{Windows!history}

Microsoft Windows is the most widely used operating system in the world. Here's how it evolved:

% --- Figure 6.2: Windows Version Timeline ---
\begin{figure}[H]
\centering
\begin{tikzpicture}[scale=0.85]
  % Timeline line
  \draw[NavyBlue, line width=2pt] (0,0) -- (15,0);
  
  % Timeline nodes
  \foreach \x/\ver/\yr/\col/\feat in {
    0.5/XP/2001/FunSky/{Start menu redesign},
    3.0/Vista/2007/FunMint/{Visual effects},
    5.5/7/2009/FunPurple/{Taskbar improved},
    8.0/8/2012/FunCoral/{Touch \& tiles},
    10.5/10/2015/ModuleBlue/{Cortana \& Edge},
    13.5/11/2021/LabGreen/{Centred taskbar}
  }{
    \node[timenode=\col] at (\x, 0) {};
    \node[above=0.4cm, font=\small\bfseries\sffamily, text=\col] at (\x, 0) {Win \ver};
    \node[below=0.3cm, font=\tiny\sffamily, text=NavyBlue] at (\x, 0) {\yr};
    \node[below=0.7cm, font=\tiny\sffamily\itshape, text=NavyBlue!70, text width=2cm, align=center] 
      at (\x, 0) {\feat};
  }
\end{tikzpicture}
\caption{A timeline of major Windows versions}
\label{fig:windows-timeline}
\end{figure}

% ---------------------------------------------------------------
\section{Functions of an Operating System}
\label{sec:os-functions}
\index{operating system!functions}

% --- Table 6.1: OS Functions Summary ---
\begin{table}[H]
\centering
\caption{Key functions of an operating system}
\label{tab:os-functions}
\renewcommand{\arraystretch}{1.3}
\begin{tabularx}{\textwidth}{c l X l}
\toprule
\rowcolor{HeaderRow}
 & \textcolor{white}{\textbf{Function}} & \textcolor{white}{\textbf{What the OS Does}} & \textcolor{white}{\textbf{Example}} \\
\midrule
\rowcolor{RowLight}
\faMicrochip & Memory & Manages RAM allocation & Opens Chrome \\
\rowcolor{RowDark}
\faHdd & Storage & Reads/writes files & Saves a document \\
\rowcolor{RowLight}
\faDesktop & Interface & Shows you the desktop & Icons, Start Menu \\
\rowcolor{RowDark}
\faNetworkWired & Network & Connects to the Internet & Wi-Fi management \\
\rowcolor{RowLight}
\faLock & Security & Protects against threats & Windows Defender \\
\bottomrule
\end{tabularx}
\end{table}

\bytetip[fun]{There are many operating systems besides Windows! \textbf{macOS} runs on Apple MacBooks, \textbf{Linux} is free and open-source (used by hackers and scientists), and \textbf{Android} runs on most smartphones. They all do the same basic job but look and feel different.}

\begin{funfact}
The name ``Windows'' comes from the fact that the operating system displays information in rectangular areas on screen called \textbf{windows}. Before Windows, people used MS-DOS, where you had to type commands in plain text --- no mouse, no icons, no colours! Imagine using a computer with only a blank screen and a blinking cursor. Scary, right?
\end{funfact}

% ---------------------------------------------------------------

\begin{reviewqs}
  \item What is an operating system? Write a definition in your own words.
  \item Name three popular operating systems.
  \item List four functions of an operating system.
  \item What was the operating system used before Windows?
  \item Why is the OS called a ``bridge'' between hardware and software?
\end{reviewqs}

\begin{challenge}
Research the history of \textbf{one operating system} (Windows, macOS, Linux, or Android). Write a short paragraph about when it was first released, who created it, and one interesting fact about it. Present your findings to the class!
\end{challenge}
